\noindent\textbf{Scope and source order.}
This chapter covers all 36 records in the supplied 1980--81 manifest and restores three omitted records, 510--512: the archive catalog dates them 23 February 1991, while the autograph P000960 reads 23 February 1981. The resulting 39 records comprise 38 talks or interventions and one catalog year divider. They are ordered by first appearance in the manuscript, including Gelfand's introductions before Beilinson and the Gelfand--Cherednik sequence on P000960. Overlapping catalog page ranges have been replaced by observed footprints. P000916 is a preceding-year continuation already indexed as 1979-488.1--.7; it is accounted for without being counted again (\citealp{mccme2026gelfand}).

\noindent\textbf{What the year contains.}
The 246 mathematical occurrences range from classroom exercises to announcements. Autumn's intersection homology, Segre families and regular holonomic systems lead into Beilinson's sustained exposition of localization, transfer bimodules and derived direct image. The chronology matters: an announcement circulating before a talk is distinguished from a later full proof, and a later geometric implementation does not automatically make an old formula new. The categorical affinity question is followed from its November introduction to its answer in the continuation.

Novikov's December lecture combines already published Pauli zero modes with the topological discussion appearing in a paper received the following day. The physical line bundle has degree \(N\); the reciprocal-torus zero-mode bundle has rank \(N\) and Chern number of absolute value one. The notebook conflates these numbers. MacPherson's Grassmannian constructions lead to positive torus orbits, harmonic forms and convexity; Khovanskii and Kazarnovskii connect Newton polyhedra, toric geometry and zero counting. Nondegeneracy, regularity and normalization conditions are restored explicitly.

Spring brings Yang--Baxter equations, commuting differential operators, algebraic geometry and \(K\)-theory, Hamiltonian kinetic theory, eigenvalue inequalities, gas dynamics and Lie semigroups. The transfer-matrix argument needs the correct auxiliary trace legs; the Weyl-operator examples need noncommutative calculation; the kinetic constructions need their actual functional and coefficient classes. The classification records the fate of the qualified statement and separately assesses the literal notebook claim.

\begin{center}
\begin{tabular}{lrrrrrr}
\toprule
Class & A & B & C & D & E & F\\
\midrule
Occurrences & 42 & 13 & 1 & 0 & 165 & 25\\
\bottomrule
\end{tabular}
\end{center}
A denotes a subsequent speaker publication; B a later development by others; C a partly developed precise question; D an undeveloped identified question; E prior knowledge; F an insufficiently identified statement. There are 16 problem identities, 94 bibliography entries and 90 correction rows. Counts describe occurrences, so a theorem repeated in several talks is counted several times. No D classification is justified by the search evidence.

\noindent\textbf{The bounded development and its limits.}
The remaining class-C item is Gelfand's conjecture that every complex algebraic curve's function field embeds in the first Weyl division ring. All genus-one and genus-two fields are positive cases: in particular Oganesyan's 2016 preprint, published in 2018, realizes every standard genus-two quintic by rational-coefficient commuting operators (\citealp[Theorem~3]{oganesyan2016rational}). The genus-two item is therefore B. The development proves an elementary obstruction for a prescribed positive-genus affine model with a nonconstant unit, supplies a complete exact certificate for a known genus-two family, and proves the injectivity step from operators to a function field. These are bounded clarifications, not a solution of the universal conjecture or a novelty claim for the known family. Independent checks caught both a notebook formal-expression error and a scalar normalization in the checked Mironov preprint.

\noindent\textbf{Continuity and evidence.}
The unextended Gelfand--Kirillov conjecture recurs under its existing identity GP-1967-21; its general form is false, with surviving simple-type cases stated cautiously. The preceding year's collective-field remarks connect to the Vlasov and hydrodynamic material without being equated with its later precise brackets. The field-embedding question is kept separate from automorphism and different-rank Weyl-field problems. The report and research ledgers identify full-text, metadata-only and unresolved evidence; a verified bibliographic link alone is never presented as a verified proof.
